% Part: incompleteness
% Chapter: arithmetization-syntax
% Section: proofs-in-lk

\documentclass[../../../include/open-logic-section]{subfiles}

\begin{document}

\olfileid{inc}{art}{plk}
\olsection{\usetoken{P}{derivation} ing $\Log{LK}$}

\begin{explain}
Supaya bisa ngaritmetisasi !!{derivation}s, kita kudu makili
!!{derivation}s minangka bilangan. Amarga !!{derivation}s iku wit
saka sekuen lan saben inferensi uga nduweni label, cara kang paling
cetha yaiku representasi rekursif: kita makili !!{derivation}
minangka tupla kang komponene yaiku sekuen pungkasan, label,
lan representasi sub-!!{derivation}s kang tumuju marang premis
inferensi pungkasan.
\end{explain}

\begin{defn}
Yen $\Gamma$ iku urutan winates saka !!{sentence}s, $\Gamma =
\tuple{!A_1, \dots, !A_n}$, mula $\Gn{\Gamma} = \tuple{\Gn{!A_1},
  \dots, \Gn{!A_n}}$.

Yen $\Gamma \Sequent \Delta$ iku sekuen, bilangan G\"odel saka
$\Gamma \Sequent \Delta$ yaiku
\[
\Gn{\Gamma \Sequent \Delta} = \tuple{\Gn{\Gamma}, \Gn{\Delta}}
\]

Yen $\pi$ iku !!a{derivation} ing $\Log{LK}$, $\Gn{\pi}$
ditemtokake kaya mangkene:
\begin{enumerate}
\item Yen $\pi$ mung dumadi saka sekuen wiwitan
  $\Gamma \Sequent \Delta$, $\Gn{\pi}$ yaiku
  \[
    \tuple{0, \Gn{\Gamma \Sequent \Delta}}.
  \]
\item Yen $\pi$ dipungkasi inferensi kanthi siji utawa loro premis,
  nduweni $\Gamma \Sequent \Delta$ minangka konklusi, lan $\pi_1$
  lan $\pi_2$ iku subbukti langsung kang pungkasane premis inferensi
  pungkasan, $\Gn{\pi}$ yaiku
  \begin{align*}
    & \tuple{1, \Gn{\pi_1}, \Gn{\Gamma \Sequent \Delta}, k} \text{ utawa}\\ 
    & \tuple{2, \Gn{\pi_1}, \Gn{\pi_2}, \Gn{\Gamma \Sequent \Delta}, k}, 
  \end{align*}
  miturut urutane, lan $k$ ditemtokake dening tabel iki miturut
  aturan kang digunakake ing inferensi pungkasan:

  \begin{tabular}{lccccccc}
    \text{Aturan:} & \LeftR{\Weakening} & \RightR{\Weakening} &
    \LeftR{\Contraction} & \RightR{\Contraction} &
    \LeftR{\Exchange} & \RightR{\Exchange} \\
    $k$: & 1 & 2 & 3 & 4 & 5 & 6 \\[2ex]
    \text{Aturan:} & \LeftR{\lnot} & \RightR{\lnot} &
    \LeftR{\land} & \RightR{\land} & 
    \LeftR{\lor} & \RightR{\lor} \\
    $k$: & 7 & 8 & 9 & 10 & 11 & 12 \\[2ex]
    \text{Aturan:} & \LeftR{\lif} & \RightR{\lif} &
    \LeftR{\lforall} & \RightR{\lforall} &
    \LeftR{\lexists} & \RightR{\lexists} \\
    $k$: & 13 & 14 & 15 & 16 & 17 & 18 \\[2ex]
    \text{Aturan:} & \Cut & = \\
    $k$: & 19 & 20
  \end{tabular}
\end{enumerate}
\end{defn}

\begin{ex}
  Gatekna !!{derivation} kang prasaja iki
  \begin{prooftree}
    \Axiom$!A \fCenter !A$
    \RightLabel{\LeftR{\land}}
    \UnaryInf$!A \land !B \fCenter !A$
    \RightLabel{\RightR{\lif}}
    \UnaryInf$\fCenter (!A \land !B) \lif !A$
  \end{prooftree}
  Bilangan G\"odel saka sekuen wiwitan yaiku $p_0 = \tuple{0,
    \Gn{!A \Sequent !A}}$. Bilangan G\"odel saka !!{derivation}
    kang pungkasane konklusi~$\LeftR{\land}$ yaiku $p_1 =
    \tuple{1, p_0, \Gn{!A \land !B \Sequent !A}, 9}$ ($1$ amarga
    $\LeftR{\land}$ nduweni siji premis, banjur bilangan G\"odel
    saka konklusi~$!A \land !B \Sequent !A$, lan $9$ iku bilangan
    kang ngodeake $\LeftR{\land}$). Mula bilangan G\"odel saka
    sakabehe !!{derivation} yaiku
    $\tuple{1, p_1, \Gn{\Sequent (!A \land !B) \lif !A}, 14}$, yaiku
  \[
  \tuple{1, \tuple{1, \tuple{0, \Gn{!A \Sequent !A}}, \Gn{!A \land !B \Sequent !A}, 9},
    \Gn{\Sequent (!A \land !B) \lif !A}, 14}.
  \]
% OLPL-170: Rong tanda kurung tutup keluwihan ing kode sekuen sumber didandani ing kene.
\end{ex}

\begin{explain}
Sawise nemtokake representasi !!{derivation}s, kita uga kudu
nuduhake yen !!{derivation}s iki bisa diolah kanthi rekursif
primitif, lan sifat sarta relasi pokoke bisa dinyatakake kanthi
cara kang padha. Sawetara operasi prasaja: upamane, yen diwenehi
bilangan G\"odel~$p$ saka !!a{derivation},
$\fn{EndSequent}(p) = (p)_{(p)_0+1}$ menehi bilangan G\"odel
saka sekuen pungkasane, lan $\fn{LastRule}(p) = (p)_{(p)_0+2}$
menehi kode aturan pungkasane. Sifat $\fn{Sequent}(s)$ kang
ditetepake dening
% OLPL-171: Jeneng EndSeq ing sumber dicocogake karo EndSequent kang digunakake sabanjure.
\[
  \len{s} = 2 \land \bforall{i<\len{(s)_0} + \len{(s)_1}}{\fn{Sent}(((s)_0 \concat (s)_1)_i)}
\]
bener kanggo $s$ yen lan mung yen $s$ iku bilangan G\"odel saka
sekuen kang dumadi saka !!{sentence}s. Sawetara sifat liyane luwih
angel. Kita bakal menehi sketsa carane. Tujuane yaiku nuduhake
yen relasi ``$\pi$ iku !!a{derivation} saka~$!A$ adhedhasar~$\Gamma$''
iku relasi rekursif primitif saka bilangan G\"odel $\pi$ lan~$!A$.
\end{explain}

\begin{prop}\ollabel{prop:followsby}
  Sifat $\fn{Correct}(p)$, kang bener yen lan mung yen inferensi
  pungkasan ing !!{derivation}~$\pi$ kanthi bilangan G\"odel~$p$
  iku bener, yaiku rekursif primitif.
\end{prop}

\begin{proof}
  $\Gamma \Sequent \Delta$ iku sekuen wiwitan yen ana
  !!a{sentence}~$!A$ saengga $\Gamma \Sequent \Delta$ padha karo
  $!A \Sequent !A$, utawa ana suku~$t$ saengga $\Gamma \Sequent
  \Delta$ padha karo $\emptyset \Sequent \eq[t][t]$.
  Ing babagan bilangan G\"odel, $\fn{InitSeq}(s)$ bener yen lan mung yen
  \begin{align*}
  \bexists{x < s}{} (\fn{Sent}(x) & \land
  s = \tuple{\tuple{x},\tuple{x}})
  \lor {}\\
  \bexists{t<s}{} (\fn{Term}(t) & \land
  s = \tuple{0, \tuple{\Gn{{\eq}(} \concat t \concat \Gn{,} \concat t \concat \Gn{)}}}).
  \end{align*}

  Kita uga kudu nuduhake yen kanggo saben aturan inferensi~$R$,
  relasi $\fn{FollowsBy}_R(p)$ iku rekursif primitif.
  $\fn{FollowsBy}_R(p)$ bener yen lan mung yen $p$ iku bilangan
  G\"odel saka !!{derivation}~$\pi$, lan sekuen pungkasane~$\pi$
  diduweni kanthi ngetrapake~$R$ kanthi bener marang
  sub-!!{derivation}s langsung saka~$\pi$.

  Kasus kang prasaja yaiku aturan \RightR{\land}. Yen $\pi$
  dipungkasi inferensi $\RightR{\land}$ kang bener, wujude kaya iki:
  \begin{prooftree}
    \AxiomC{}
    \RightLabel{$\pi_1$}
    \Deduce$\Gamma \fCenter \Delta, !A$
    
    \AxiomC{}
    \RightLabel{$\pi_2$}
    \Deduce$\Gamma \fCenter \Delta, !B$

    \RightLabel{\RightR\land}
    \BinaryInf$\Gamma \fCenter \Delta, !A \land !B$
  \end{prooftree}
  Mula inferensi pungkasan ing !!{derivation} $\pi$ iku
  penerapan $\RightR{\land}$ kang bener yen lan mung yen ana
  urutan !!{sentence}s $\Gamma$ lan $\Delta$, uga loro
  !!{sentence}s $!A$ lan~$!B$, saengga sekuen pungkasane $\pi_1$
  yaiku $\Gamma \Sequent \Delta, !A$, sekuen pungkasane $\pi_2$
  yaiku $\Gamma \Sequent \Delta, !B$, lan sekuen pungkasane $\pi$
  yaiku $\Gamma \Sequent \Delta, !A \land !B$. Iki mung perlu
  diowahi dadi bilangan G\"odel. Yen $s = \Gn{\Gamma \Sequent \Delta}$,
  mula $(s)_0 = \Gn{\Gamma}$ lan $(s)_1 = \Gn{\Delta}$.
  Mula $\fn{FollowsBy}_{\RightR{\land}}(p)$ bener yen lan mung yen
\begin{align*}
& \bexists{g < p}{\bexists{d < p}{\bexists{a < p}{\bexists{b < p}{\quad}}}} \\
& \qquad \fn{EndSequent}(p) = 
  \tuple{g, d \concat 
    \tuple{\Gn{(} \concat a \concat \Gn{\land} \concat b \concat \Gn{)}}} \land {} \\
& \qquad \fn{EndSequent}((p)_1) = \tuple{g, d \concat \tuple{a}} \land {}\\
& \qquad \fn{EndSequent}((p)_2) = \tuple{g, d \concat \tuple{b}} \land {}\\
& \qquad (p)_0 = 2 \land \fn{LastRule}(p) = 10.
\end{align*}
Baris-baris kasebut, miturut urutane, nyatakake ``ana
urutan~($\Gamma$) kanthi bilangan G\"odel~$g$, ana
urutan~($\Delta$) kanthi bilangan G\"odel~$d$, !!a{formula}~($!A$)
kanthi bilangan G\"odel~$a$, lan !!a{formula}~($!B$) kanthi
bilangan G\"odel~$b$,'' saengga ``sekuen pungkasane $\pi$ yaiku
$\Gamma \Sequent \Delta, !A \land !B$,'' ``sekuen pungkasane $\pi_1$
yaiku $\Gamma \Sequent \Delta, !A$,'' ``sekuen pungkasane $\pi_2$
yaiku $\Gamma \Sequent \Delta, !B$,'' lan ``$\pi$ nduweni loro
subderivasi langsung lan aturan inferensi pungkasane yaiku
$\RightR\land$ (kanthi nomer~$10$).''

Inferensi pungkasan ing~$\pi$ iku penerapan $\RightR\lexists$
kang bener yen lan mung yen ana urutan $\Gamma$ lan $\Delta$,
!!a{formula}~$!A$, variabel~$x$, lan suku~$t$, saengga sekuen
pungkasane $\pi$ yaiku $\Gamma \Sequent \Delta, \lexists[x][!A]$
lan sekuen pungkasane $\pi_1$ yaiku $\Gamma \Sequent \Delta,
\Subst{!A}{t}{x}$. Dadi, ing babagan bilangan G\"odel,
$\fn{FollowsBy}_{\RightR\lexists}(p)$ bener yen lan mung yen
\begin{align*}
  & \bexists{g<p}{\bexists{d<p}{\bexists{a<p}{\bexists{x<p}{\bexists{t<p}{\quad}}}}}\\
  & \qquad \fn{EndSequent}(p) = 
  \tuple{
    g,
    d \concat \tuple{\Gn{\lexists} \concat x \concat a}
  } \land {}\\
  & \qquad \fn{EndSequent}((p)_1) = 
  \tuple{
    g,
    d \concat \tuple{\fn{Subst}(a, t, x)}
  } \land {}\\
  & \qquad (p)_0 = 1 \land \fn{LastRule}(p) = 18.
\end{align*}
Sabanjure kita netepake $\fn{Correct}(p)$ minangka
\begin{multline*}
  \fn{Sequent}(\fn{EndSequent}(p)) \land {}\\ 
  [(\fn{LastRule}(p) = 1 \land 
  \fn{FollowsBy}_{\LeftR\Weakening}(p)) \lor \dots \lor {}\\
  (\fn{LastRule}(p) = 20 \land \fn{FollowsBy}_{\eq}(p)) \lor {}\\
  (p)_0 = 0 \land \fn{InitSeq}(\fn{EndSequent}(p))]
\end{multline*}
% OLPL-172: InitialSeq ing sumber ora padha karo InitSeq kang wis ditetepake.
Baris kapisan mesthekake yen sekuen pungkasane~$p$ pancen
sekuen kang dumadi saka !!{sentence}s. Baris pungkasan nyakup
kasus nalika $p$ mung sekuen wiwitan.
% OLPL-173: Variabel d ing panjelasan sumber diganti p supaya cocog karo Correct(p).
\end{proof}

\begin{prob}
  Tetepna sifat-sifat iki kaya ing
  \olref[inc][art][plk]{prop:followsby}:
  \begin{enumerate}
  \item $\fn{FollowsBy}_{\Cut}(p)$,
  \item $\fn{FollowsBy}_{\LeftR\lif}(p)$,
  \item $\fn{FollowsBy}_{\eq}(p)$,
  \item $\fn{FollowsBy}_{\RightR{\lforall}}(p)$.
  \end{enumerate}
  Kanggo sifat pungkasan, kowé uga kudu nuduhake yen kanthi
  rekursif primitif bisa dipriksa apa inferensi pungkasan saka
  !!{derivation} kanthi bilangan G\"odel~$p$ nyukupi syarat
  eigenvariabel, yaiku eigenvariabel~$a$ saka $\RightR{\lforall}$
  ora muncul ing sekuen pungkasan.
  \end{prob}
  

\begin{prop}
  \ollabel{prop:deriv}
  Relasi $\fn{Deriv}(p)$, kang bener yen $p$ iku bilangan G\"odel
  saka !!{derivation}~$\pi$ kang bener, iku rekursif primitif.
\end{prop}

\begin{proof}
  !!^a{derivation}~$\pi$ iku bener yen saben inferensine
  ngetrapake aturan kanthi bener, yaiku yen saben
  sub-!!{derivation}s-e dipungkasi inferensi kang bener.
  Mula $\fn{Deriv}(p)$ bener yen lan mung yen
  \[
  \bforall{i<\len{\fn{SubtreeSeq}(p)}}{\fn{Correct}((\fn{SubtreeSeq}(p))_i)}.
  \]
% OLPL-174: Deriv(d) ing sumber diganti Deriv(p), lan kurung tutup Correct ditambahake.
\end{proof}


\begin{prop}
Upamane $\Gamma$ iku himpunan rekursif primitif saka
!!{sentence}s. Relasi $\Prf[\Gamma](x, y)$ kang nyatakake
``$x$ iku kode saka !!a{derivation}~$\pi$ kanggo
$\Gamma_0 \Sequent !A$ kanthi sawijining $\Gamma_0 \subseteq \Gamma$
kang winates, lan $y$ iku bilangan G\"odel saka~$!A$''
iku rekursif primitif.
\end{prop}

\begin{proof}
Upamane ``$y \in \Gamma$'' diwenehake dening predikat
rekursif primitif~$R_\Gamma(y)$. Kita kudu nuduhake yen
$\Prf[\Gamma](x, y)$, kang bener yen lan mung yen $y$ iku
bilangan G\"odel saka ukara~$!A$ lan $x$~iku kode saka
$\Log{LK}$-!!{derivation} kanthi sekuen pungkasan
$\Gamma_0 \Sequent !A$, iku rekursif primitif.

Miturut proposisi sadurunge, sifat $\fn{Deriv}(x)$, kang bener
yen lan mung yen $x$ iku kode saka !!{derivation}~$\pi$ kang
bener ing $\Log{LK}$, iku rekursif primitif. Yen $x$ iku kode
kaya mangkono, $\fn{EndSequent}(x)$ iku kode sekuen pungkasane~$\pi$,
mula $(\fn{EndSequent}(x))_0$ iku kode sisih kiwa sekuen
pungkasan lan $(\fn{EndSequent}(x))_1$ kode sisih tengene.
Mula ``sisih tengen sekuen pungkasane~$\pi$ yaiku~$!A$'' bisa
dinyatakake minangka
$\len{(\fn{EndSequent}(x))_1} = 1 \land ((\fn{EndSequent}(x))_1)_0 =
y$. Sisih kiwa sekuen pungkasane $\pi$ mesthi winates;
kita mung kudu nyatakake yen saben ukara ing kono dadi anggota
ing~$\Gamma$. Mula $\Prf[\Gamma](x, y)$ bisa ditetepake dening
% OLPL-175: Kode ukara ing katerangan sumber kudu y, cocog karo Prf(x,y) lan rumus pungkasan.
\begin{align*}
\Prf[\Gamma](x, y) \defiff {}&
\fn{Deriv}(x) \land {} \\
& \bforall{i <
  \len{(\fn{EndSequent}(x))_0}}{R_\Gamma(((\fn{EndSequent}(x))_0)_i)} \land {}\\
& \len{(\fn{EndSequent}(x))_1} = 1 \land ((\fn{EndSequent}(x))_1)_0 = y.
\end{align*}
\end{proof}

\end{document}
